\chapter{Liouville's Theorem and Magnetic Forces}

This lecture has two tightly linked halves. We begin with Liouville's theorem, not as a technical afterthought, but as a statement that Hamiltonian mechanics preserves information: phase-space flow neither converges nor diverges, and a region of phase space keeps its volume even while its shape changes. We then turn, in a deliberately different direction, to a charged particle in a magnetic field, where the new feature is a force that depends on velocity and yet still fits naturally into the Lagrangian and Hamiltonian framework. These companion notes follow Leonard Susskind's lecture, with supporting transcript and frame curation prepared with LazyingArt LLC and Video2Book.

\section{Liouville's Theorem as Information Conservation}

Susskind begins by insisting that Liouville's theorem deserves to be taken seriously before we even prove it. In Hamiltonian mechanics it is central, and in quantum mechanics its close analog is unitarity. In that sense the theorem is not merely geometric; it is a classical expression of information conservation.

The picture is simple and powerful. A phase space is the space of possible configurations of a system. In the Hamiltonian setting the natural coordinates are canonical pairs \((q_i,p_i)\), but the first intuition does not depend on this notation. Imagine that every point of phase space is occupied by a system, and let time run. Then the points move, and they move as a flow.

What kind of flow can classical mechanics allow? Each point has one future and one past. No trajectory can split into two, for then the equations of motion would cease to determine what happens next. No two distinct trajectories can merge into one, for then the present would fail to determine the past. In this sense the flow is both deterministic and reversible. Susskind phrases the deeper geometric content as follows: the flow in phase space is incompressible.

It is useful to think of phase space as filled with a uniform dust of points. Each point may be surrounded by a tiny volume element, and as the flow carries that dust along, those little volume elements are deformed, stretched, and sheared, but not squeezed away. A region may become elongated, twisted, or filamentary; it does not lose its phase-space volume.

\section{Phase-Space Flow, Friction, and the Meaning of ``No Convergence''}

At once an objection arises from ordinary experience. Friction seems to make many different initial conditions collapse onto the same final state. A block sliding on a table comes to rest. Dice tumble in many different ways, yet end on one of a few visible outcomes. Is this not convergence in phase space?

The lecture's answer is that such examples are not fundamental counterexamples; they are examples of incomplete description. If we ignore the microscopic degrees of freedom of the table, the air, or the internal structure of the system, then an effective description can indeed look dissipative. But the contraction comes from having thrown away variables.

This is exactly parallel to the familiar appearance of energy loss. A sliding object seems to lose kinetic energy, but a more careful accounting reveals that the table warms up. Energy was not destroyed; degrees of freedom were ignored. Likewise, if a phase-space volume appears to contract, the moral is not that Liouville's theorem failed, but that the system was described too coarsely.

That is why Susskind refuses to identify Liouville's theorem with mere determinism of differential equations. Determinism alone says that a point has a definite future. Liouville's theorem says more: the Hamiltonian flow preserves phase-space volume. One must not flatten the statement into uniqueness of solutions. The geometric content is stronger.

\section{From Fluid Incompressibility to Hamiltonian Incompressibility}

Before returning to phase space, the lecture makes a strategic detour. Let us first recall what incompressibility means for an ordinary flow.

In one dimension the conclusion is almost trivial. Suppose points move along a line with local velocity \(v_x(x)\). If a small interval is to keep the same density of points, then the number entering must equal the number leaving. A variation of velocity from one end of the interval to the other would cause clumping or thinning. Therefore the velocity must satisfy
\begin{equation}
\frac{\partial v_x}{\partial x}=0.
\end{equation}
So an incompressible one-dimensional flow is just rigid translation: every point moves with the same instantaneous velocity.

Two dimensions are more interesting. Consider a small rectangle of sides \(\Delta x\) and \(\Delta y\). The net change in the number of points inside the rectangle comes from flux through its four sides. The vertical sides contribute a term proportional to
\[
\left(\frac{\partial v_x}{\partial x}\right)\Delta x\,\Delta y,
\]
and the horizontal sides contribute
\[
\left(\frac{\partial v_y}{\partial y}\right)\Delta x\,\Delta y.
\]
For an incompressible flow the sum must vanish, so we obtain
\begin{equation}
\frac{\partial v_x}{\partial x}+\frac{\partial v_y}{\partial y}=0.
\end{equation}
In general dimensions this becomes
\begin{equation}
\nabla\cdot \mathbf v = 0,
\end{equation}
or, in index notation,
\begin{equation}
\sum_i \frac{\partial v_i}{\partial x_i}=0.
\end{equation}

Now we import this language into Hamiltonian mechanics. In phase space the coordinates are the canonical variables \(q_i\) and \(p_i\), and the local phase-space velocity is the collection of \(\dot q_i\) and \(\dot p_i\). Hamilton's equations read
\begin{equation}
\dot p_i=-\frac{\partial H}{\partial q_i},
\qquad
\dot q_i=\frac{\partial H}{\partial p_i}.
\end{equation}
Therefore the phase-space divergence is
\begin{equation}
\sum_i\left(\frac{\partial \dot q_i}{\partial q_i}+\frac{\partial \dot p_i}{\partial p_i}\right).
\end{equation}
Substituting Hamilton's equations gives
\begin{align}
\sum_i\left(\frac{\partial \dot q_i}{\partial q_i}+\frac{\partial \dot p_i}{\partial p_i}\right)
&=
\sum_i\left(
\frac{\partial}{\partial q_i}\frac{\partial H}{\partial p_i}
-
\frac{\partial}{\partial p_i}\frac{\partial H}{\partial q_i}
\right) \\
&=
\sum_i\left(
\frac{\partial^2 H}{\partial q_i\,\partial p_i}
-
\frac{\partial^2 H}{\partial p_i\,\partial q_i}
\right)=0.
\end{align}
The mixed partial derivatives cancel pairwise. This is the whole proof.

\begin{theorem}[Liouville's theorem]
Hamiltonian flow preserves phase-space volume.
\end{theorem}

\begin{proof}
The phase-space divergence vanishes. Therefore Hamiltonian flow is incompressible, and an evolving region of phase space retains its volume.
\end{proof}

What is preserved, then, is not the metric distance between points, but the volume of a region transported by the flow. That is the key point of the next stage of the lecture.

\section{Area, Not Distance, in Phase Space}

To make the geometry sharp, Susskind turns to the simplest possible mechanical system: one coordinate, one momentum, and a free-particle Lagrangian. In the lecture he sets the mass equal to \(1\), so
\begin{equation}
L=\frac{\dot x^2}{2},
\qquad
p_x=\frac{\partial L}{\partial \dot x}=\dot x.
\end{equation}

\begin{figure}
\centering
\includegraphics[width=0.58\textwidth]{lecture_07_figure_02.png}
\caption{Free-particle Lagrangian and the definition of \(P_x\).}
\end{figure}

Now change coordinates. Let
\begin{equation}
y=\alpha x,
\qquad
x=\frac{y}{\alpha},
\qquad
\dot x=\frac{\dot y}{\alpha}.
\end{equation}
The same Lagrangian becomes
\begin{equation}
L=\frac{\dot y^2}{2\alpha^2},
\end{equation}
and the momentum conjugate to \(y\) is
\begin{equation}
p_y=\frac{\partial L}{\partial \dot y}
=
\frac{\dot y}{\alpha^2}
=
\frac{\dot x}{\alpha}
=
\frac{p_x}{\alpha}.
\end{equation}
So when the coordinate axis is stretched by \(\alpha\), the conjugate momentum axis is shrunk by the same factor. Hence a small phase-space rectangle transforms according to
\begin{equation}
\Delta y\,\Delta p_y
=
(\alpha\,\Delta x)\left(\frac{\Delta p_x}{\alpha}\right)
=
\Delta x\,\Delta p_x.
\end{equation}

This is the right invariant object. The widths and heights need not be preserved separately, but the area is preserved. That is why time evolution itself is naturally viewed as a canonical transformation: it changes coordinates on phase space in a way that preserves the phase-space area element.

\subsection{Question \& Answer}

\paragraph{Question.} How can distances between phase-space points change while phase-space area stays fixed?

\paragraph{Answer.} Because the flow can stretch in one direction while compressing in another. A square may shear into a parallelogram of the same area. Two points that were initially nearby in one direction may separate, while other pairs become closer, and yet the area of the whole patch remains unchanged. Liouville's theorem is not a statement about rigidity. It is a statement about incompressibility.

\begin{remark}
This phase-space area becomes especially meaningful in quantum mechanics. The lecture emphasizes that Planck's constant measures an area scale in phase space, not a position scale and not a momentum scale separately.
\end{remark}

\section{Chaos, Coarse-Graining, and Entropy}

The next conceptual obstacle is chaos. If a small region of phase space is evolved under a chaotic Hamiltonian, it may stretch, fold, and develop increasingly fine filaments. To the eye it appears to spread all over phase space. One might think the occupied volume has grown.

The lecture's answer is subtle and important. Fine-grained evolution preserves volume exactly. The filaments may become extraordinarily thin and elaborate, but no gap opens, no piece breaks off, and the total phase-space volume remains the same. In this exact sense, chaos is not a counterexample to Liouville's theorem. It is one of its most vivid illustrations.

At the same time, phase-space area carries probabilistic meaning. For a chaotic system, the relative probability of finding the system in two small regions is tied to the corresponding phase-space areas, if probability is understood operationally as the fraction of time spent in each region. This is why the invariant phase-space measure matters so much in statistical mechanics.

\subsection{Question \& Answer}

\paragraph{Question.} If chaotic motion spreads out, why does Liouville's theorem not fail?

\paragraph{Answer.} Because the spreading is compensated by compression in transverse directions. A small patch may become long and thin, or highly folded, but its exact volume is unchanged. What grows is not the fine-grained phase-space volume, but the coarse-grained volume we must use when our resolution is finite.

This brings us to coarse-graining. Instead of following exact phase points, suppose we only resolve finite cells. Then a fine filamentary structure eventually threads through many such cells, and the union of those resolvable cells occupies a much larger apparent region than the original compact patch. In this coarse-grained sense the effective phase-space volume grows.

That is where entropy enters. In the lecture it is described heuristically but very clearly: entropy measures the detectable phase-space volume occupied by the system, and more precisely it is the logarithm of that volume. Exact Hamiltonian dynamics preserves information. Entropy increase appears because our description no longer follows the microscopic filaments with sufficient precision.

\section{Why Magnetic Fields Are the Right Next Example}

After this long geometric discussion, Susskind turns sharply to a new example: a charged particle moving in an electromagnetic field. The reason for choosing magnetism is precise. Up to now the forces under discussion depended on position but not on velocity. A magnetic force changes that.

For a particle of charge \(q\) moving with velocity \(\mathbf v\) in a magnetic field \(\mathbf B\), the force is
\begin{equation}
\mathbf F = q\,\mathbf v\times \mathbf B.
\end{equation}
This is a velocity-dependent force. Friction is also velocity dependent, but the similarity is only superficial. Friction is not part of fundamental microscopic mechanics; it is an effective force that appears after degrees of freedom have been discarded. Magnetic forces, by contrast, do come from an action principle. They admit a Lagrangian formulation, a Hamiltonian formulation, and a conserved energy, provided we treat the magnetic field as an externally specified, time-independent background.

To formulate the mechanics, however, one does not work directly with \(\mathbf B\). One introduces a vector potential \(\mathbf A\) such that
\begin{equation}
\mathbf B = \nabla\times \mathbf A.
\end{equation}
The component formulas needed later are
\begin{equation}
B_x=\frac{\partial A_z}{\partial y}-\frac{\partial A_y}{\partial z},
\qquad
B_y=\frac{\partial A_x}{\partial z}-\frac{\partial A_z}{\partial x},
\qquad
B_z=\frac{\partial A_y}{\partial x}-\frac{\partial A_x}{\partial y}.
\end{equation}

\subsection{Question \& Answer}

\paragraph{Question.} Why can we not write the mechanics directly in terms of \(\mathbf B\) alone?

\paragraph{Answer.} Because Newton's equation \(m\mathbf a=q\,\mathbf v\times \mathbf B\) is not by itself a Lagrangian or Hamiltonian formulation. To write an action, to define canonical momenta, and to construct the Hamiltonian, one needs an object whose derivatives reproduce the magnetic field. That object is the vector potential \(\mathbf A\). The equations of motion ultimately depend only on \(\mathbf B\), but the variational formalism is built from \(\mathbf A\).

\section{Charged Particle in a Magnetic Field: Action, Momentum, and Energy}

The lecture now proceeds exactly as an experienced theorist often proceeds: make the simplest plausible guess for the action, and then check that it works. The free-particle term must remain,
\[
\int \frac12 mv^2\,dt.
\]
The new ingredient is a line integral of the vector potential along the particle's path. Thus the action is taken to be
\begin{equation}
S=\int \frac12 mv^2\,dt + q\int A_i\,dx_i.
\end{equation}
The board crop in the lecture shows only the beginning of this second term, but the full line-integral form is the standard completion consistent with the spoken derivation.

\begin{figure}
\centering
\includegraphics[width=0.62\textwidth]{lecture_07_figure_03.png}
\caption{Action guess with kinetic term and vector-potential coupling.}
\end{figure}

Since \(dx_i=\dot x_i\,dt\), the action can be written in the usual velocity form,
\begin{equation}
S=\int \left(\frac12 m\dot x_i^2 + qA_i\dot x_i\right)dt,
\end{equation}
and therefore the Lagrangian is
\begin{equation}
L=\frac12 m\dot x_i^2 + qA_i\dot x_i.
\end{equation}

The first new surprise comes immediately. The canonical momentum is not just \(m\dot x_i\), but
\begin{equation}
p_i=\frac{\partial L}{\partial \dot x_i}=m\dot x_i+qA_i.
\end{equation}
In particular,
\begin{equation}
p_x=m\dot x+qA_x,
\qquad
p_y=m\dot y+qA_y,
\qquad
p_z=m\dot z+qA_z.
\end{equation}
The lecture calls \(m\dot x_i\) the mechanical momentum, while \(p_i\) is the canonical momentum. These differ by \(qA_i\).

Now we verify that the Lagrange equations reproduce the Lorentz force. Let us work out the \(z\)-component. The Euler--Lagrange equation is
\begin{equation}
\frac{d}{dt}\frac{\partial L}{\partial \dot z}=\frac{\partial L}{\partial z}.
\end{equation}
The left-hand side is
\begin{equation}
\frac{d}{dt}\frac{\partial L}{\partial \dot z}
=
\frac{d}{dt}(m\dot z+qA_z)
=
m\ddot z+q\dot A_z,
\end{equation}
while the right-hand side is
\begin{equation}
\frac{\partial L}{\partial z}
=
q\dot x\,\frac{\partial A_x}{\partial z}
+
q\dot y\,\frac{\partial A_y}{\partial z}
+
q\dot z\,\frac{\partial A_z}{\partial z}.
\end{equation}

\begin{figure}
\centering
\includegraphics[width=0.62\textwidth]{lecture_07_figure_04.png}
\caption{First term in \(\partial L/\partial z\) from the vector potential.}
\end{figure}

Because the particle moves through space, even a time-independent magnetic field produces a nonzero \(\dot A_z\) along the trajectory:
\begin{equation}
\dot A_z
=
\frac{\partial A_z}{\partial x}\dot x
+
\frac{\partial A_z}{\partial y}\dot y
+
\frac{\partial A_z}{\partial z}\dot z.
\end{equation}
Substituting this into the Euler--Lagrange equation gives
\begin{align}
m\ddot z
&=
q\dot x\left(\frac{\partial A_x}{\partial z}-\frac{\partial A_z}{\partial x}\right)
+
q\dot y\left(\frac{\partial A_y}{\partial z}-\frac{\partial A_z}{\partial y}\right).
\end{align}
Now use
\[
B_y=\frac{\partial A_x}{\partial z}-\frac{\partial A_z}{\partial x},
\qquad
B_x=\frac{\partial A_z}{\partial y}-\frac{\partial A_y}{\partial z},
\]
so that
\begin{equation}
m\ddot z=q(v_xB_y-v_yB_x)=q(\mathbf v\times \mathbf B)_z.
\end{equation}
The \(x\)- and \(y\)-components follow by cyclic permutation. Thus the guessed action really does reproduce
\[
m\mathbf a=q\,\mathbf v\times \mathbf B.
\]

At this stage the lecture makes a second conceptual turn and asks about energy. Assume that \(\mathbf A\), and hence \(\mathbf B\), has no explicit time dependence. Then the Hamiltonian will be conserved. By definition,
\begin{equation}
H=\sum_i p_i\dot x_i-L.
\end{equation}
If we calculate it first in terms of velocities, the result is strikingly simple:
\begin{align}
H
&=
(m\dot x_i+qA_i)\dot x_i-\left(\frac12 m\dot x_i^2+qA_i\dot x_i\right) \\
&=
\frac12 m\dot x_i^2
=
\frac12 mv^2.
\end{align}
The magnetic field contributes nothing explicitly to the energy in velocity variables.

To use Hamilton's equations, however, we must express the Hamiltonian in terms of canonical momenta. From
\begin{equation}
p_i=m\dot x_i+qA_i
\end{equation}
we solve for the velocity:
\begin{equation}
\dot x_i=\frac{p_i-qA_i}{m}.
\end{equation}
Therefore
\begin{equation}
H=\frac{1}{2m}\sum_i (p_i-qA_i)^2
=
\frac{1}{2m}\bigl(\mathbf p-q\mathbf A\bigr)^2.
\end{equation}
This is the form appropriate to Hamiltonian mechanics.

\subsection{Question \& Answer}

\paragraph{Question.} Why does the magnetic field alter the canonical momentum but not the energy?

\paragraph{Answer.} Because the magnetic field changes the relation between canonical variables and mechanical motion. The mechanical momentum is
\[
\mathbf p_{\mathrm{mech}}=m\mathbf v=\mathbf p-q\mathbf A,
\]
so the canonical momentum differs from the mechanical one by \(q\mathbf A\). But the energy depends on speed, and magnetic forces are perpendicular to the velocity. They deflect the trajectory without changing \(|\mathbf v|\). That is why
\[
H=\frac12 mv^2
\]
in velocity variables, even though in canonical variables the Hamiltonian becomes
\[
H=\frac{1}{2m}(\mathbf p-q\mathbf A)^2.
\]

The lecture closes by observing that one can now take this canonical Hamiltonian and verify, as an exercise, that Hamilton's equations reproduce the same Lorentz-force equations of motion. The formalism is self-consistent: one and the same physics appears once in Newton's language, once in the Lagrangian language, and once in the Hamiltonian language.

\section{Summary}

Liouville's theorem tells us that Hamiltonian mechanics is an incompressible flow on phase space. A region may shear, stretch, and fold, but its volume is preserved. That is the classical content of information conservation. It leads naturally to the invariant area element of phase space, to canonical transformations, and to the statistical meaning of phase-space measure.

The second half of the lecture shows that magnetic forces, though velocity dependent, do not lie outside this framework. They require a vector potential, shift canonical momentum away from mechanical momentum, and yet preserve the simple kinetic form of the energy because magnetic fields do no work. In that sense the lecture moves from the geometry of Hamiltonian flow to a concrete dynamical example, and the two parts fit together exactly as they should.